\documentclass[11pt]{article}
\usepackage[T1]{fontenc}
\usepackage{lmodern,microtype}
\usepackage[margin=1in]{geometry}
\usepackage{amsmath,amssymb,amsthm,mathtools}
\usepackage{graphicx,booktabs,tabularx,array}
\usepackage[numbers,sort&compress]{natbib}
\usepackage{xcolor}
\usepackage[colorlinks=true,linkcolor=blue!45!black,citecolor=blue!45!black,%
            urlcolor=blue!45!black]{hyperref}
\usepackage{enumitem}

\theoremstyle{plain}
\newtheorem{theorem}{Theorem}
\newtheorem{proposition}[theorem]{Proposition}
\newtheorem{lemma}[theorem]{Lemma}
\newtheorem{corollary}[theorem]{Corollary}
\theoremstyle{definition}
\newtheorem{definition}[theorem]{Definition}
\theoremstyle{remark}
\newtheorem{remark}[theorem]{Remark}

\newcommand{\E}{\mathbb E}
\newcommand{\Prob}{\mathbb P}
\newcommand{\eps}{\varepsilon}
\newcommand{\Cov}{\operatorname{Cov}}
\newcommand{\Var}{\operatorname{Var}}
\newcommand{\TV}{\operatorname{TV}}
\newcommand{\kl}{\operatorname{kl}}
\newcommand{\sgn}{\operatorname{sign}}
\DeclareMathOperator{\Bin}{Bin}
\newcommand{\cp}{c_{\mathrm p}}
\setlength{\emergencystretch}{2em}

\title{When single-lineage measurements fail to determine treatment order:\\
sister dependence, branching extinction, and decision-focused assays}
\author{}
\date{}

\begin{document}
\maketitle

\begin{abstract}
Measurements sufficient to describe a retained cell lineage need not determine
which intervention minimises the survival of the whole family. We construct a
finite-type branching source in which complete retained-branch histories have
identical laws across a family of sister couplings, while two interventions
applied for exactly the same durations have opposite preferred orders. An exact
finite-duration formula identifies the decision boundary and shows that it is
strictly positive, not at zero covariance. Two new ingredients then control
recurring descendant division: the extinction probability is monotone in the
descendant division rate, and a sensitivity reserve replaces the crude
occupation bound. Together they certify the reversal at $h=\log\frac87$ for
descendant division rates up to $\eps=1/30$ instead of $1/500$, and they
certify it for the \emph{entire} range $0\le\eps\le1$ at pulses of length
$\log\frac{50}{49}$ rather than $10^{-6}$ --- four orders of magnitude of
pulse length --- so that the phenomenon survives descendants dividing exactly
as fast as the founder. A sister-agreement measurement repairs the missing
information: forty independent division records give $95\%$ correct choice on a
separated constant-treatment class, now throughout $\eps\le1/4$, and exact
binomial inversion gives full-class rules with an explicit unresolved output.
We distinguish calibration needed for a shifted decision threshold from
calibration-free sign information, price correlated read errors as an additive
covariance reserve, and express residual uncertainty in excess-survival-risk
units. The source is synthetic: its rate ratios, instantaneous readout and
independence assumptions are stated explicitly, and the conclusion is a
statement about an observation protocol, not about a clinically effective
regimen.
\end{abstract}

\tableofcontents

\section{Introduction}

\subsection{Synopsis}

A cell's descendants form a family, but a measurement may follow only one
branch of that family. Two sources can produce exactly the same observed branch
while coupling the states of sister cells differently. We show that this hidden
dependence can reverse the preferred order of two specified treatment phases,
so that no amount of data from the admitted protocol resolves the comparison. A
paired sister measurement recovers the information the decision actually needs,
and exact binomial calculations quantify when it supports a declaration. The
examples are synthetic and demonstrate an information requirement, not a
clinically effective regimen.

\subsection{Inference about a cell and a decision about a family}

Inferring proliferation and phenotype transitions, and choosing an
intervention, are related but different tasks. Branching inference from cell
numbers and fractions can exploit information well beyond mean trajectories
\citep{GunnarssonFooLeder2023}, and optimal-dosing analysis studies
interventions under a specified source and response law
\citep{GunnarssonMagnussonFoo2025}. Here we ask a prior question: does an
explicitly declared observation experiment determine a particular intervention
comparison, when the source is not fully known?

The distinction matters at division. A retained daughter can have a completely
known state distribution while its dependence with the omitted sister is
unknown. Family extinction, however, uses \emph{products} of daughter
continuation probabilities, so it is sensitive to exactly that dependence. A
schedule comparison can therefore turn on information erased by the observation
protocol. This is not an assertion that all cell counts or all lineage trees are
uninformative. We prove equality of one specific complete pruned-record law,
including its timestamps, event labels, death and censoring records, and we say
precisely which protocols are excluded.

\subsection{What this paper establishes}

The paper assembles a single source-specific package.

\begin{enumerate}[nosep,leftmargin=*]
\item \textbf{An exact observation equivalence} (Theorem~\ref{thm:branch}).
      Two bounded-rate finite-type sources with common founder law, life-event
      rates, observation channels and one-daughter transition marginals, but
      different sister couplings, produce identical complete retained-branch
      records on any bounded horizon, under every shared policy measurable from
      the retained history.
\item \textbf{An exact finite-duration treatment-order boundary}
      (Proposition~\ref{prop:boundary}). For the constructed source the
      difference of extinction probabilities between the chronological
      schedules $AB$ and $BA$ factors completely; the boundary
      $\cp(e^{-h})$ is strictly positive and strictly increasing in the pulse
      length, with $\cp(e^{-h})=7h/8+O(h^2)$, and it leaves the feasible
      covariance range at a single point.
\item \textbf{Monotonicity in the descendant division rate}
      (Lemma~\ref{lem:mono}). Family extinction is nonincreasing in the
      descendant division rate $\eps$, under a supersolution condition that
      Lemma~\ref{lem:criteria} verifies for the source at hand and, at every
      pulse length where the reversal exists, for an arbitrary descendant
      offspring law. This is new here and it halves every perturbation reserve
      at once, because the two schedules can only move in the same direction.
\item \textbf{A sensitivity reserve} (Proposition~\ref{prop:reserve}). A
      differential-inequality bound replaces the coupling bound ``probability
      that an extra division occurs'' by a bound that also charges how much
      such a division can matter. It is between eight and thirty-one times
      smaller than the occupation reserve it replaces, at no cost in
      generality: it holds for every $\eps\ge0$.
\item \textbf{A much larger certified range for the reversal}
      (Theorem~\ref{thm:pulse}, Theorem~\ref{thm:short},
      Table~\ref{tab:improve}). At $h=\log\frac87$ the opposite preferences at
      $c=\mp1/5$ now persist through $\eps=1/30$, and at $c=\mp1/4$ through
      $\eps=1/16$. At the shorter pulse $h=\log\frac{50}{49}$ the reversal
      holds for \emph{every} $\eps\in[0,1]$, that is, with descendants dividing
      exactly as fast as the founder.
\item \textbf{A uniform decision band} (Corollary~\ref{cor:band}). For every
      pulse length, every $\eps$ and every $c$, the preferred order is
      determined outside an explicit band of half-width proportional to $\eps$
      and, for small $h$, to about $11h$. The previously available statement of
      this shape had half-width $16000\,h$ and applied only at $\eps\le1$.
\item \textbf{Finite-sample decision rules for the repaired measurement}
      (Theorem~\ref{thm:forty}, Proposition~\ref{prop:abstain}), with an
      unresolved output, exact binomial validity over the whole admitted source
      class, and calculated resolution probabilities.
\item \textbf{Calibration and error-model accounting}
      (\S\ref{sec:calibration}), including a calibration-free sign rule that is
      now usable at finite pulse length, and an additive reserve
      (Proposition~\ref{prop:corr}) that prices correlated read errors instead
      of assuming them away.
\end{enumerate}

Items 3--6 and 8 are new relative to the earlier treatment of this source;
items 1, 2 and 7 are retained, with the numerical content recomputed.
Table~\ref{tab:improve} summarises the quantitative change.

\begin{table}[tb]
\centering\small
\caption{What the new certificates change. ``Previous'' refers to the
occupation-coupling certificate; every entry is for the same source, the same
schedules and the same exact arithmetic. Factors are rounded down.}
\label{tab:improve}
\begin{tabularx}{\linewidth}{@{}Xllr@{}}
\toprule
Statement & Previous & This paper & Factor\\\midrule
Reversal at $c=\mp\tfrac15$, $h=\log\frac87$ & $\eps\le1/500$ & $\eps\le1/30$ & $16.8$\\
Reversal at $c=\mp\tfrac14$, $h=\log\frac87$ & --- & $\eps\le1/16$ & ---\\
Reversal throughout $0\le\eps\le1$ & $h\le10^{-6}$ & $h=\log\frac{50}{49}\approx.0202$ & $>2\times10^4$\\
Half-width of the undecided band in $c$ & $16000\,h$ & $\approx11\,h\,\eps$ & $>10^3$\\
Separated $40$-division class & $\eps\le.001$ & $\eps\le1/4$ & $250$\\
Full-class pulse rule at $.95$ resolution & $m=1600$, $\eps\le.001$ & $m=800$, $\eps\le.01$ & $2$ and $10$\\
Matched direct-endpoint budget & $7\,814\,272$ roots & $5\,377\,920$ roots & $1.45$\\
\bottomrule
\end{tabularx}
\end{table}

\subsection{Sign conventions and a map of the decision tasks}
\label{sec:conventions}

Throughout, $F_a$ is the probability that the family is \emph{extinct} by the
deadline under schedule $a$, and $R_a=1-F_a$ is the survival risk that the
decision minimises. We write
\[
 D:=F_{AB}-F_{BA}=-(R_{AB}-R_{BA}).
\]
\emph{A positive $D$ favours $AB$; a positive $R_{AB}-R_{BA}$ favours $BA$.}
The same convention is used for constant actions, with $D_C:=F_A-F_B$, so a
positive $D_C$ favours $A$. All reported ``gaps'' are in units of absolute
survival risk. The hidden coordinate is the founder sister covariance $c$; for
binary state indicators with marginal $1/2$ the correlation is $4c$, so
$c=\pm1/5$ means correlation $\pm0.8$ and $c=\pm1/4$ means correlation $\pm1$.
These are strong sister associations, not perturbations of independence.

Four decision tasks appear, and they must not borrow one another's parameter
ranges. Table~\ref{tab:map} is the map.

\begin{table}[tb]
\centering\small
\caption{The four decision tasks, their thresholds and their certified
descendant-division ranges. ``Low $c$'' and ``high $c$'' refer to the hidden
founder covariance. $h=T=\log\frac87$ unless stated.}
\label{tab:map}
\begin{tabularx}{\linewidth}{@{}lXll@{}}
\toprule
Task & Threshold on $c$ and which side wins & Certified $\eps$ & Where\\\midrule
Pulses $AB$ vs $BA$ & $\cp\approx.13706$; low $c\Rightarrow AB$ & $\le1/30$ ($c=\mp\frac15$) & Thm~\ref{thm:pulse}\\
 & & $\le1/16$ ($c=\mp\frac14$) & Thm~\ref{thm:pulse}\\
Short pulses, $h=\log\frac{50}{49}$ & $\cp\approx.01813$; low $c\Rightarrow AB$ & $\le1$ ($c=\mp\frac14$) & Thm~\ref{thm:short}\\
Constant $A$ vs $B$, separated & $c_*\approx.08432$; high $c\Rightarrow A$ & $\le1/4$ & Thm~\ref{thm:forty}\\
Either, full class with abstention & own threshold & $\le1/100$ & Prop~\ref{prop:abstain}\\
\bottomrule
\end{tabularx}
\end{table}

\subsection{Related work, and what is not claimed}

Decision-focused experimental design, prediction identifiability, partial
identification and minimax regret are established subjects
\citep{Flassig2015,Manski2004,Stoye2012}. Information lower bounds for
identification from repeated sampling are standard
\citep{KaufmannCappeGarivier2016,Tsybakov2009}, and we use them as background
rather than as new inequalities. Multi-type branching processes and their
backward equations are classical
\citep{Harris1963,AthreyaNey1972,KimmelAxelrod2015}; the comparison arguments in
\S\ref{sec:recurring} use the standard theory of quasimonotone differential
inequalities \citep{Walter1970,Walter1998,Smith1995}.

Our contribution is a source-specific connection: an exact intervention-order
boundary hidden from retained lineages, sharp control of the perturbation that
transfers it to recurring division, and a paired measurement with quantitative
decision guarantees. We make no exhaustive priority claim for the combined
statement, and we do not claim to have established a clinically effective
treatment, a generally optimal control policy, a universal failure of lineage
tracing, or an optimal laboratory sample budget. Section~\ref{sec:scope}
states the boundaries in detail.

\section{The source, the endpoint and the admitted record}\label{sec:source}

\subsection{The branching source}

There are three labelled types, $P,S,T$, and initially one $P$. The founder
divides at rate one, without death or switching, and is replaced by an ordered
pair of daughters with law
\begin{equation}\label{eq:kernel}
 K_c=\begin{pmatrix}1/4+c&1/4-c\\1/4-c&1/4+c\end{pmatrix},
 \qquad -\tfrac14\le c\le\tfrac14 .
\end{equation}
The row and column index the two daughter states, not mother types. An $S$ or
$T$ cell divides at rate $\eps\ge0$, using the same joint law, and dies at the
action-dependent rates
\begin{equation}\label{eq:rates}
 (d_S^A,d_T^A)=(2,4),\qquad (d_S^B,d_T^B)=(3,3).
\end{equation}
Different cells have independent event clocks and independent kernel draws
conditional on their realised genealogy. Their future subtrees are independent
\emph{conditional on both daughter states}; those states themselves need not be
independent. Replacing the joint offspring probabilities by products of
marginals would delete the mechanism under study.

The source \eqref{eq:kernel} has sister dependence but no marginal
mother-to-daughter phenotype memory: every daughter marginal is $(1/2,1/2)$.
That separation is deliberate, and \S\ref{sec:refinement} proves that the
conclusions survive a refinement with genuine marginal persistence. The
response law \eqref{eq:rates}, the preparation and the remaining source
assumptions are known; $c$ is the hidden coordinate, and we sometimes allow
$\eps$ to remain unknown inside a certified interval. This is not
extrapolation to an unmeasured action whose response is unrestricted.

\subsection{Schedules and the endpoint}

For a deterministic schedule $a$ and horizon $T$ the survival risk is
\begin{equation}\label{eq:risk}
 R_a(c,\eps)=\Prob^a_P(N_T>0),\qquad F_a=1-R_a .
\end{equation}
The principal comparison is chronological $AB$ against $BA$, each phase of
length $h$, so $T=2h$. These two schedules use exactly the same interventions
for exactly the same durations and differ only in order; there is no exposure
asymmetry to explain the difference away. Constant $A$ against constant $B$ at a
common horizon supplies a second, easier decision task. Their equal sums of
death hazards, $2+4=3+3$, are a mathematical actuator normalisation and are not
an assertion that the two actions have equal dose, toxicity or resource cost.

\subsection{The admitted observation protocol}

One observation protocol follows a fresh founder until it divides or until an
independent deadline of rate $\lambda>0$. It records the elapsed time, which
event occurred and, at a division, the state or marker of a uniformly and
independently chosen daughter. The record then ends. Its densities are
\begin{equation}\label{eq:record}
 \lambda e^{-(1+\lambda)t}\ \ \text{(deadline)},\qquad
 \tfrac12 e^{-(1+\lambda)t}\ \ \text{(division, either retained state)} .
\end{equation}
A fixed emission channel acts on the state marginal. All of these densities are
independent of $c$ and of $\eps$. Independent founders, not repeated
observations of one lineage, are the independent sample units.

The richer protocol treated in Theorem~\ref{thm:branch} continues along the
retained branch indefinitely: at every division the observer retains one
daughter by a fresh uniform choice independent of the states, discards the
sister and her entire subtree, and continues. Figure~\ref{fig:protocol} shows
this record and the paired repair studied in \S\ref{sec:measurement}.

\begin{figure}[tb]
\centering\includegraphics[width=\linewidth]{figures/protocol.pdf}
\caption{The admitted observation and its repair. In the admitted record the
unretained sister and all of her descendants are absent, even though the
record is otherwise complete: states, event labels, timestamps, death and
censoring are all retained, for as long as the horizon allows. A paired read at
the founder's division supplies one family-level feature --- whether the two
daughters agree --- and that single bit is what the schedule comparison needs.
No selection on subsequent survival is permitted in either protocol.}
\label{fig:protocol}
\end{figure}

\subsection{Why the backward equation is the source probability}
\label{sec:backward}

Bounded per-cell division rates give domination by a Yule process of rate
$\max(1,\eps)$, so the chronological source is nonexplosive: the independent
Yule holding-time expectations have divergent harmonic sum and summable centred
variances, hence their sum diverges almost surely. The dominating population
has finite integrated cell time on every bounded horizon, which also excludes
accumulation of death events.

Define the offspring polynomial
\begin{equation}\label{eq:G}
 G_c(s,t)=(\tfrac14+c)(s^2+t^2)+(\tfrac12-2c)\,st
 =\tfrac14(s+t)^2+c\,(s-t)^2 .
\end{equation}
Conditioning on the first event of a cell yields, for a constant action $a$,
the field
\begin{equation}\label{eq:field}
 \Phi_a^\eps(p,s,t)=\begin{pmatrix}
 G_c(s,t)-p\\
 d_S^a(1-s)+\eps\,(G_c(s,t)-s)\\
 d_T^a(1-t)+\eps\,(G_c(s,t)-t)\end{pmatrix}.
\end{equation}
The product structure in $G_c$ is legitimate only after conditioning on both
daughter states. The first-event integral equation on finite trees, followed by
bounded convergence under nonexplosion, gives $v'=\Phi_a^\eps(v)$, $v(0)=0$,
where the independent variable $u$ is \emph{time remaining to the deadline}.
The field is polynomial, hence locally Lipschitz, and the unit cube is
invariant because the field points inward on every face; uniqueness therefore
identifies the solution with the actual extinction probabilities. For
chronological phases, conditioning at the phase boundaries expresses the
extinction vector as the flow composition $F^{t_1}_{a_1}\circ\cdots\circ
F^{t_k}_{a_k}(0)$, in which the phases are evaluated in reverse chronological
order. No population-size truncation is involved.

We write $A(u)=e^{-R_S(u)}$ and $B(u)=e^{-R_T(u)}$, where $R_i(u)$ is the death
hazard of a type-$i$ cell integrated over the last $u$ units of the schedule.
At $\eps=0$ the daughter coordinates decouple and
\begin{equation}\label{eq:s0t0}
 s^0(u)=1-A(u),\qquad t^0(u)=1-B(u).
\end{equation}

\section{Identical retained lineages, different family decisions}

\begin{theorem}[Retained-branch equivalence]\label{thm:branch}
Consider two bounded-rate finite-type branching sources with the same founder
law, the same life-event rates, the same observation channels and the same
one-daughter transition marginals under each action, but possibly different
sister couplings. At every observed division retain one daughter by a fresh
uniform choice independent of the states, discard the sister and her subtree,
and continue only along the retained branch. Record every observed state, life
event and time, including death and censoring. Then the complete records on a
bounded horizon have identical laws. The conclusion remains true for shared
action and stopping policies measurable from the retained history, and for
shared external randomisation.
\end{theorem}

\begin{proof}
Given the retained history, the next event clock and the event type have the
same conditional law in both sources, because they depend only on the retained
state and the currently chosen action. At a division the retained state is
distributed according to the common daughter marginal, and the readout channel
is common. Induction over events proves equality of the law of every
finite-event record, including the terminal stopping or death flag. Equal
histories induce equal actions and equal stopping decisions under a shared
policy, and common external randomisation can be coupled identically.

Along a single branch the total event intensity is bounded by a constant, since
the action menu is finite. The number of events by a bounded horizon is
therefore dominated by a Poisson count, so truncating the record after $n$
events has a tail probability tending to zero uniformly in both sources.
Passing to the limit proves equality of the complete history laws.
\end{proof}

\begin{remark}
The theorem excludes: records of hidden family outcomes, records of both
daughter states at any division, survivor-biased sampling, and any rule that
chooses a daughter using information about the omitted sister. It establishes
much more than equality of means, and it is stable under post-processing:
because the two record laws are equal, \emph{no} measurable statistic,
reconstruction algorithm or learned model applied to those records can
distinguish the sources. Collecting more time points on the pruned branch, or
more independent pruned branches, cannot help either.
\end{remark}

The algebraic reason for the phenomenon is the identity
\begin{equation}\label{eq:square}
 G_{c+u}(s,t)-G_c(s,t)=u\,(s-t)^2 ,
\end{equation}
immediate from the second form in \eqref{eq:G}. A change of $c$ at fixed
marginals is invisible to a retained branch, but it changes a family outcome
whenever the two daughter continuation probabilities differ. The quantity the
decision depends on is a product of continuation risks, not a marginal state
frequency. A schedule comparison then asks whether the two interventions weight
the hidden direction $(s-t)^2$ differently.

\section{An exact finite-duration treatment-order boundary}\label{sec:boundary}

Set $x=e^{-h}$. For a daughter whose death rate is $r$ in the first phase and
$s$ in the second, put
\begin{equation}\label{eq:B}
 B_{r,s}(x)=\frac{x^{s+1}-x^{r+s}}{r-1}+\frac{x^{2}-x^{s+1}}{s-1},
 \qquad r,s\ne1 .
\end{equation}
Indeed a daughter born at time $t<h$ survives to $2h$ with probability
$e^{-r(h-t)-sh}$, while one born at $h\le t<2h$ survives with probability
$e^{-s(2h-t)}$; multiplying by the founder division density $e^{-t}$ and
integrating separately over the two intervals gives \eqref{eq:B}.

At $\eps=0$ the two daughter subtrees are independent given their states, so we
may expand the joint extinction product conditional on the founder's division
time. For first-phase rates $(r_1,r_2)$ and second-phase rates $(s_1,s_2)$,
\begin{align}\label{eq:two-phase}
 F={}&1-x^2-B_{r_1,s_1}-B_{r_2,s_2}\notag\\
 &+(\tfrac14+c)\bigl(B_{2r_1,2s_1}+B_{2r_2,2s_2}\bigr)
 +(\tfrac12-2c)\,B_{r_1+r_2,\,s_1+s_2} ,
\end{align}
the founder that never divides surviving and contributing no extinction mass.
Substituting $(2,4;3,3)$ and $(3,3;2,4)$ and collecting powers of $x$ gives the
complete factorisation
\begin{equation}\label{eq:D}
 D:=F_{AB}-F_{BA}=\frac{x^{2}(x-1)^{3}}{210}\bigl\{4c\,\mathcal A(x)+(x-1)\mathcal Q(x)\bigr\},
\end{equation}
where
\begin{align*}
 \mathcal A(x)&=(x^4+x^3+x^2+x+1)(3x^5+6x^4+9x^3+12x^2+8x+4),\\
 \mathcal Q(x)&=3x^8+12x^7+30x^6+60x^5+98x^4+137x^3+170x^2+159x+66 .
\end{align*}
Equivalently, writing
\begin{equation}\label{eq:slope}
 s_0(x)=\frac{4\mathcal A(x)\,x^2(1-x)^3}{210}>0,
\end{equation}
we have the affine form $D=s_0(x)\,(\cp(x)-c)$, with the boundary defined next.

\begin{proposition}[Exact phase boundary]\label{prop:boundary}
For $0<x<1$ and $\eps=0$, the schedule $AB$ has strictly smaller survival risk
if and only if $c<\cp(x)$, where
\begin{equation}\label{eq:cp}
 \cp(x)=\frac{(1-x)\mathcal Q(x)}{4\mathcal A(x)} .
\end{equation}
This threshold is strictly positive on $(0,1)$ and strictly increasing in the
pulse length $h$. There is exactly one $x_0\in(.80120,.80121)$ with
$\cp(x_0)=1/4$. For $x<x_0$ the schedule $AB$ is strictly preferred throughout
the feasible range $-1/4\le c\le1/4$; for $x>x_0$ that range contains both
preferences. Moreover $\cp(e^{-h})=\tfrac78h+O(h^2)$.
\end{proposition}

\begin{proof}
$\mathcal A$ and $\mathcal Q$ have positive coefficients and the prefactor
$x^2(x-1)^3/210$ in \eqref{eq:D} is negative on $(0,1)$; hence $D>0$, that is
$R_{AB}<R_{BA}$, exactly when $4c\mathcal A+(x-1)\mathcal Q<0$, i.e.\ below the
threshold \eqref{eq:cp}. Direct differentiation gives
$\cp'(x)=-P(x)/\bigl(4\mathcal A(x)^2\bigr)$ with
\begin{align}\label{eq:P}
 P(x)={}&735x^{10}+4410x^{9}+11655x^{8}+20580x^{7}+28560x^{6}\notag\\
 &+31290x^{5}+26740x^{4}+17850x^{3}+9030x^{2}+3080x+420 ,
\end{align}
all of whose coefficients are positive; so $\cp$ decreases strictly in $x$ and
therefore increases strictly in $h$. Its endpoint limits are $33/8$ at $x=0$
and $0$ at $x=1$, so by continuity and strict monotonicity there is exactly one
$x_0$ with $\cp(x_0)=1/4$; rational substitution at $.80120$ and $.80121$
isolates it, and monotonicity --- not a floating-point root count --- proves
uniqueness. Finally $\mathcal A(1)=210$ and $\mathcal Q(1)=735$, so
$\cp(e^{-h})=(735/840)h+O(h^2)=\tfrac78h+O(h^2)$.
\end{proof}

\begin{remark}\label{rem:positive}
The finite-duration decision threshold is \emph{not} at $c=0$. This is the
single most consequential correction to the naive picture: the leading-order
calculation of \S\ref{sec:legacy} sees only the cubic term $-4ch^3$ and
therefore places the boundary at zero covariance, whereas at $h=\log\frac87$
the boundary sits at $\cp=8967219299/65423564404\approx.1370640591$, more than
half way to the feasible ceiling. Consequently a measurement that determines
only the \emph{sign} of $c$ does not determine the preferred order at finite
pulse length; see \S\ref{sec:calibration}. Proposition~\ref{prop:boundary} is
also a positive identification result of independent interest: for $x<x_0$ the
hidden parameter remains completely unknown and yet its entire feasible range
gives the same decision. What matters is whether the decision boundary meets
the observationally equivalent set, not whether every parameter is recovered.
\end{remark}

At $x=7/8$, that is $h=\log\frac87\approx.133531$, the exact values are
\begin{equation}\label{eq:witnesses}
\begin{aligned}
 \cp&=\tfrac{8967219299}{65423564404}\approx.1370640591, &
 s_0&=\tfrac{114491237707}{32985348833280}\approx3.47097\times10^{-3},\\
 D(-\tfrac14)&=\tfrac{5635}{4194304}\approx1.34349\times10^{-3}, &
 D(\tfrac14)&=-\tfrac{25860351307}{65970697666560}\approx-3.91998\times10^{-4},\\
 D(-\tfrac15)&=\tfrac{771817626293}{659706976665600}\approx1.16994\times10^{-3}, &
 D(\tfrac15)&=-\tfrac{48037425121}{219902325555200}\approx-2.18449\times10^{-4}.
\end{aligned}
\end{equation}
The contrast has opposite signs at $c=-1/5$ and $c=+1/5$, and again at
$c=\mp1/4$: negative covariance favours $AB$ and positive covariance favours
$BA$. Figure~\ref{fig:boundary} displays the boundary and the contrast.

\begin{figure}[tb]
\centering\includegraphics[width=\linewidth]{figures/boundary.pdf}
\caption{Left: the exact analytic threshold $\cp(e^{-h})$ and its first-order
approximation $7h/8$, with the feasible ceiling $c=1/4$ and the unique crossing
$h_0=-\log x_0$. Below $h_0$ the schedule $AB$ wins for every feasible
covariance; above it both preferences occur, and the gap between the exact
curve and $7h/8$ is what a sign-only measurement cannot see. Right: the exact
reference contrast at $h=\log\frac87$, with the two witnesses marked, enclosed
by the proved sensitivity band at $\eps\le10^{-2}$ (Proposition~\ref{prop:reserve})
and, for comparison, by the occupation band at the ten times smaller
$\eps\le10^{-3}$. Both bands are analytic enclosures, not solver-tolerance
estimates. Negative contrast favours $BA$.}
\label{fig:boundary}
\end{figure}

\section{Recurring descendant division}\label{sec:recurring}

The reference calculation of \S\ref{sec:boundary} is exact only at $\eps=0$.
This section transfers it to $\eps>0$, and it is where the present treatment
differs most from what was previously available. We prove two structural facts
--- a monotonicity and a sensitivity bound --- and only then evaluate constants.

Throughout this section $a$ is a schedule with phases of constant death rates,
$U$ is its total horizon, and $R_S,R_T$ are the integrated death hazards over
the last $u$ units, as in \S\ref{sec:backward}.

\subsection{Extinction is monotone in the descendant division rate}

We allow the descendants to use an offspring law $G$ that may differ from the
founder's $G_c$; this costs nothing and is what makes \S\ref{sec:refinement}
work. Thus $G$ is any polynomial $w_{SS}s^2+w_{ST}st+w_{TT}t^2$ with
nonnegative coefficients summing to one, and \eqref{eq:field} is read with
$G_c$ in the first coordinate and $G$ in the other two. Everything below is
driven by one inequality on the \emph{reference} trajectory.

\begin{lemma}[Monotonicity]\label{lem:mono}
Assume $d_S\le d_T$ in every phase, and assume the supersolution condition
\begin{equation}\label{eq:supersol}
 G\bigl(s^0(u),t^0(u)\bigr)\ \le\ s^0(u)\qquad\text{for all }u\in[0,U].
\end{equation}
Then for every $\eps\ge0$ the extinction vectors satisfy
$v^{\eps}(u)\le v^{0}(u)$ componentwise on $[0,U]$. In particular
$F^{\eps}_a\le F^{0}_a$: recurring descendant division can only \emph{reduce}
family extinction.
\end{lemma}

\begin{proof}
Since $d_S\le d_T$ we have $s^0\le t^0$, so \eqref{eq:supersol} says
$G(s^0,t^0)\le\min(s^0,t^0)$ and both $\eps$-terms of \eqref{eq:field} are
$\le0$ at $v^0$. Hence $(v^0)'=\Phi^0_a(v^0)\ge\Phi^{\eps}_a(v^0)$: the
reference trajectory is a supersolution of the $\eps$-system.

On $[0,1]^3$ each component of $\Phi^{\eps}_a$ is nondecreasing in the other
two coordinates, because $\partial_sG_c=2(\tfrac14+c)s+(\tfrac12-2c)t\ge0$ and
$\partial_tG_c\ge0$ there (both coefficients $\tfrac14+c$ and $\tfrac12-2c$ are
nonnegative for $|c|\le\tfrac14$), because the same holds for $G$, and because
$\eps\ge0$. The system is therefore quasimonotone.

On the first phase $v^{\eps}$ solves $v'=\Phi^{\eps}_a(v)$ with $v(0)=0$ and
$v^0$ is a supersolution of the same system with the same initial value; both
remain in the invariant cube $[0,1]^3$. The comparison theorem for
quasimonotone systems \citep[\S10]{Walter1970} (see also
\citep[Ch.~3]{Smith1995}, \citep[\S{}III.12]{Walter1998}) gives
$v^{\eps}\le v^0$ on that phase. At a phase boundary the ordering is inherited
as initial data, \eqref{eq:supersol} holds on the next phase as well, and the
argument repeats.
\end{proof}

Condition \eqref{eq:supersol} says that dividing now is never better for
extinction than remaining a single cell of the more robust type. Two criteria
cover every use we make of it, one specific to the founder's kernel and
uniform in the pulse length, the other valid for an arbitrary descendant
kernel on a restricted range of pulse lengths.

\begin{lemma}[Criteria for \eqref{eq:supersol}]\label{lem:criteria}
\hfill
\begin{enumerate}[nosep,leftmargin=*,label=(\alph*)]
\item If $G=G_c$ with $|c|\le\tfrac14$ and if $d_S\le d_T\le2d_S$ in every
      phase, then \eqref{eq:supersol} holds, for every pulse length.
\item For an \emph{arbitrary} descendant kernel $G$, \eqref{eq:supersol} holds
      as soon as $\bigl(t^0(u)\bigr)^2\le s^0(u)$ for all $u$. For the two
      actions \eqref{eq:rates} and each of the schedules $AB$, $BA$, $A$, $B$
      with common phase length $h$, this in turn holds if and only if
      \begin{equation}\label{eq:xflat}
       2x^2-x^6\ \ge\ 1,\qquad\text{that is}\qquad
       x\ \ge\ x_\flat:=\Bigl(\tfrac{\sqrt5-1}{2}\Bigr)^{1/2}=.7861513\ldots
      \end{equation}
\end{enumerate}
Since $x_\flat<x_0$, criterion (b) is available at \emph{every} pulse length at
which both preferences occur in the feasible covariance range.
\end{lemma}

\begin{proof}
(a) By the second form in \eqref{eq:G} and $|c|\le\tfrac14$,
$G_c(s,t)=\tfrac14(s+t)^2+c(s-t)^2\le\tfrac12(s^2+t^2)$. The rate hypothesis
integrates to $R_S\le R_T\le2R_S$, so with $A,B$ as in \eqref{eq:s0t0} we get
$0<B\le A\le1$ and $A^2=e^{-2R_S}\le e^{-R_T}=B$; hence $A^2+B^2\le B+B=2B$ and
\[
 \tfrac12\bigl((1-A)^2+(1-B)^2\bigr)=1-(A+B)+\tfrac12(A^2+B^2)\le1-A=s^0 .
\]

(b) Every monomial of $G$ is a product of two of $s^0,t^0$, each at most
$\max(s^0,t^0)=t^0$, and the coefficients sum to one, so
$G(s^0,t^0)\le (t^0)^2$. The condition $(t^0)^2\le s^0$ reads
$(1-B)^2\le1-A$, i.e.
\begin{equation}\label{eq:ab-crit}
 A\ \le\ B\,(2-B).
\end{equation}
We check \eqref{eq:ab-crit} on each phase, writing $x=e^{-h}$.
For a phase with $d_S=d_T$ one has $A=B$ and \eqref{eq:ab-crit} reads
$A\le A(2-A)$, i.e.\ $A\le1$, which always holds; this settles the schedule
$B$, the first backward phase of $AB$ and the second backward phase of $BA$
except for the accumulated prefactors, treated next.
For the first backward phase of $BA$, and for the constant schedule $A$, the
rates are $(2,4)$ from the deadline, so $A=e^{-2u}$ and $B=A^2$ with
$A\in[x^2,1]$; \eqref{eq:ab-crit} becomes $2A-A^3\ge1$. The cubic
$A^3-2A+1=(A-1)(A^2+A-1)$ shows $2A-A^3\ge1$ exactly for
$A\ge\tfrac{\sqrt5-1}{2}$, and the smallest value of $A$ on the phase is
$x^2$, so the condition is $x^2\ge\tfrac{\sqrt5-1}{2}$, i.e.\ \eqref{eq:xflat}.
For the second backward phase of $BA$ we have $A=x^2z$, $B=x^4z$ with
$z=e^{-3(u-h)}\in[x^3,1]$, and \eqref{eq:ab-crit} becomes
$2x^2-x^6z\ge1$, whose worst case $z=1$ is again \eqref{eq:xflat}.
For the second backward phase of $AB$ we have $A=x^3w^2$, $B=x^3w^4$ with
$w=e^{-(u-h)}\in[x,1]$, and \eqref{eq:ab-crit} becomes $f(w):=2w^2-x^3w^6\ge1$;
since $f'(w)=2w\,(2-3x^3w^4)$ vanishes at most once on $(0,1]$ and $f$ is there
first increasing, the minimum of $f$ over $[x,1]$ is attained at an endpoint,
and $f(1)=2-x^3\ge1$ while $f(x)=2x^2-x^9\ge2x^2-x^6\ge1$ by \eqref{eq:xflat}.
Conversely the case $A=e^{-2u}$ above shows that \eqref{eq:xflat} is necessary.
Finally $x_\flat^2=\tfrac{\sqrt5-1}{2}<.61804<.64192<x_0^2$, so
$x_\flat<x_0$.
\end{proof}

\begin{remark}
Criterion (a) applies to our source with $(d^A_S,d^A_T)=(2,4)$ saturating
$d_T\le2d_S$ exactly, and it is uniform in $h$; criterion (b) is sharp in $x$
but indifferent to the descendant kernel. Both are sufficient conditions for
\eqref{eq:supersol}, which is what the analysis actually uses.
\end{remark}

\subsection{A sensitivity reserve}

Lemma~\ref{lem:mono} fixes the sign of the perturbation. The next proposition
bounds its size, and does so more sharply than the obvious coupling argument.

\begin{definition}[Two reserves]\label{def:reserves}
For $i\in\{S,T\}$ put
\begin{equation}\label{eq:sigmaL}
 L_i(u)=e^{-R_i(u)}\!\int_0^u\! e^{R_i(\tau)}\,d\tau ,
 \qquad
 \sigma_i(u)=e^{-R_i(u)}\!\int_0^u\!\bigl(e^{R_i(\tau)}-1\bigr)d\tau ,
\end{equation}
and define the \emph{occupation reserve} and the \emph{sensitivity reserve}
\begin{equation}\label{eq:VW}
 V_a=\int_0^{U}\! e^{-(U-u)}\bigl[L_S(u)+L_T(u)\bigr]du ,
 \qquad
 W_a=\int_0^{U}\! e^{-(U-u)}\,t^0(u)\bigl[\sigma_S(u)+\sigma_T(u)\bigr]du .
\end{equation}
\end{definition}

Here $L_i(u)$ is the expected time a type-$i$ daughter present at remaining
time $u$ stays alive, so $V_a$ is the expected integrated live-daughter time of
the reference process --- the quantity a naive coupling produces. In
$\sigma_i$ that live time is weighted by $s_i^0(\tau)=1-e^{-R_i(\tau)}$, the
probability that the daughter's own subtree would have died anyway, and in
$W_a$ the result is weighted again by $t^0(u)\le1$. Both weights are at most
one, so $W_a\le V_a$; the point of the proposition is that the weights are in
fact small.

\begin{proposition}[Sensitivity reserve]\label{prop:reserve}
Assume $d_S\le d_T$ in every phase together with the supersolution condition
\eqref{eq:supersol}. Then for every $c\in[-\tfrac14,\tfrac14]$ and every
$\eps\ge0$,
\begin{equation}\label{eq:reserve}
 0\ \le\ F^0_a-F^{\eps}_a\ \le\ \eps\,W_a\ \le\ \eps\,V_a .
\end{equation}
Consequently, for two schedules $a,b$ and $D^{\eps}=F^{\eps}_a-F^{\eps}_b$,
\begin{equation}\label{eq:contrast-band}
 D^{0}-\eps W_a\ \le\ D^{\eps}\ \le\ D^{0}+\eps W_b .
\end{equation}
\end{proposition}

\begin{proof}
Let $\Delta=v^0-v^{\eps}$, which is nonnegative by Lemma~\ref{lem:mono}.
Subtracting the two systems componentwise gives the exact identities
\begin{align}
 \Delta_s'&=-d_S\Delta_s+\eps\bigl(s^{\eps}-G(s^{\eps},t^{\eps})\bigr),
 \label{eq:ds}\\
 \Delta_t'&=-d_T\Delta_t+\eps\bigl(t^{\eps}-G(s^{\eps},t^{\eps})\bigr),
 \nonumber\\
 \Delta_p'&=G_c(s^0,t^0)-G_c(s^{\eps},t^{\eps})-\Delta_p .
 \label{eq:dp}
\end{align}
Since $G\ge0$ on the cube and $s^{\eps}\le s^0$, the forcing in \eqref{eq:ds}
is at most $\eps\,s^0$. Multiplying by the integrating factor $e^{R_S}$ and
using $\Delta_s(0)=0$,
\[
 \Delta_s(u)\ \le\ \eps\int_0^u e^{-(R_S(u)-R_S(\tau))}s^0(\tau)\,d\tau
 \ =\ \eps\,\sigma_S(u),
\]
the last equality because $e^{R_S(\tau)}s^0(\tau)=e^{R_S(\tau)}-1$. Identically
$\Delta_t(u)\le\eps\,\sigma_T(u)$.

For \eqref{eq:dp}, the segment joining $(s^{\eps},t^{\eps})$ to $(s^0,t^0)$
lies in $[0,s^0]\times[0,t^0]$, so by the mean value theorem there is a point
$(\tilde s,\tilde t)$ on it with
$G_c(s^0,t^0)-G_c(s^{\eps},t^{\eps})
 =\partial_sG_c(\tilde s,\tilde t)\Delta_s+\partial_tG_c(\tilde s,\tilde t)\Delta_t$.
Because $\tilde s\le s^0\le t^0$ and $\tilde t\le t^0$, and because
$2(\tfrac14+c)+(\tfrac12-2c)=1$,
\[
 \partial_sG_c(\tilde s,\tilde t)=2(\tfrac14+c)\tilde s+(\tfrac12-2c)\tilde t\le t^0 ,
\]
and the same bound holds for $\partial_tG_c$. Hence
$\Delta_p'\le t^0(\Delta_s+\Delta_t)-\Delta_p$, and with $\Delta_p(0)=0$,
\[
 F^0_a-F^{\eps}_a=\Delta_p(U)\le\int_0^{U}e^{-(U-u)}t^0(u)\bigl(\Delta_s+\Delta_t\bigr)(u)\,du
 \le\eps W_a .
\]
The comparison $W_a\le V_a$ follows from $t^0\le1$ and $\sigma_i\le L_i$. For
\eqref{eq:contrast-band}, subtract the two one-sided bounds and use that both
deficits are nonnegative: $D^{\eps}=D^0-(F^0_a-F^{\eps}_a)+(F^0_b-F^{\eps}_b)$
with the first deficit in $[0,\eps W_a]$ and the second in $[0,\eps W_b]$.
\end{proof}

\begin{remark}\label{rem:twoimprovements}
Two distinct improvements are combined in \eqref{eq:contrast-band}. First,
Lemma~\ref{lem:mono} makes the band \emph{one-sided in each direction}: without
knowing the sign, one would have to write $|D^\eps-D^0|\le\eps(W_a+W_b)$,
losing a factor close to two because $W_{AB}\approx W_{BA}$. Second,
$W_a$ itself is far below $V_a$. The two gains multiply, and neither costs any
generality: \eqref{eq:reserve} is valid for \emph{all} $\eps\ge0$, not only for
small $\eps$.
\end{remark}

Both reserves are elementary integrals of exponentials with integer rates, so
for rational $x=e^{-h}$ each evaluates to an element of
$\mathbb Q+\mathbb Q\,h$. Appendix~\ref{app:reserve} gives the closed forms;
Table~\ref{tab:reserves} gives the values at $x=7/8$.

\begin{table}[tb]
\centering\small
\caption{The two reserves at $x=7/8$. The occupation values reproduce the
previously published constants exactly, which validates the evaluation; the
sensitivity values are the ones used from here on. Here $h=\log\frac87$, and
the decimal column is an enclosure computed from a rational bracket for $h$ of
width below $10^{-30}$.}
\label{tab:reserves}
\begin{tabularx}{\linewidth}{@{}llXr@{}}
\toprule
Schedule & Occupation $V_a$ & Sensitivity $W_a$ & $V_a/W_a$\\\midrule
$AB$ & $\frac{957367}{18874368}\approx.0507231$
     & $-\frac{6624078645151}{54975581388800}+\frac{10392084486973}{10995116277760}h\approx.00571652$ & $8.87$\\[2pt]
$BA$ & $\frac{2552741}{50331648}\approx.0507184$
     & $-\frac{520588799010593}{3958241859993600}+\frac{6795050968489}{6597069766656}h\approx.00601852$ & $8.43$\\[2pt]
$A$  & $\approx.0149943$
     & $-\frac{11540424829}{105696460800}+\frac{206906567}{251658240}h\approx.000601292$ & $24.9$\\[2pt]
$B$  & $\approx.0149740$
     & $-\frac{79873}{1228800}+\frac{321391}{655360}h\approx.000483482$ & $31.0$\\
\bottomrule
\end{tabularx}
\end{table}

\subsection{The reversal under recurring division}

\begin{theorem}[Exposure-matched decision reversal]\label{thm:pulse}
Take the source of \S\ref{sec:source} with $h=\log\frac87$ and compare the
chronological schedules $AB$ and $BA$.
\begin{enumerate}[nosep,leftmargin=*]
\item At $c=-1/5$ the schedule $AB$ has strictly smaller survival risk, and at
      $c=+1/5$ the schedule $BA$ does, uniformly for $0\le\eps\le1/30$. The
      absolute risk gap is at least $1.78\times10^{-5}$ throughout that range.
\item At $c=-1/4$ and $c=+1/4$ the same opposite preferences hold uniformly for
      $0\le\eps\le1/16$, with absolute risk gap at least $1.58\times10^{-5}$.
\end{enumerate}
The exact suprema of the admissible ranges are
\[
 \eps_\star(\tfrac15)=\min\Bigl(\frac{D(-\tfrac15)}{W_{AB}},\ \frac{-D(\tfrac15)}{W_{BA}}\Bigr)
 \approx.0362961 ,
 \qquad
 \eps_\star(\tfrac14)\approx.0651318 .
\]
For each fixed $\eps$ the complete retained-branch laws of the two sources are
identical, under every shared admissible retained-history policy.
\end{theorem}

\begin{proof}
At $\eps=0$ the exact evaluation \eqref{eq:witnesses} of the factored contrast
\eqref{eq:D} gives $D(-1/5)>0>D(1/5)$ and $D(-1/4)>0>D(1/4)$. Here the
descendant kernel is the founder's, and $d_S\le d_T\le2d_S$ holds for both
actions, so Lemma~\ref{lem:criteria}(a) gives \eqref{eq:supersol} and
Proposition~\ref{prop:reserve} applies with $a=AB$, $b=BA$, yielding
\eqref{eq:contrast-band}. At $c=-1/5$, positivity of $D^{\eps}$ therefore
follows from $D(-1/5)>\eps W_{AB}$, and at $c=+1/5$ negativity follows from
$-D(1/5)>\eps W_{BA}$; the binding constraint is the second. Substituting the
values of Table~\ref{tab:reserves} and clearing denominators in exact rational
arithmetic proves both claims at $\eps=1/30$, with the stated gap, and
similarly at $\eps=1/16$ for $c=\mp1/4$. Theorem~\ref{thm:branch} supplies the
observation equality.
\end{proof}

\begin{remark}
The previously available certificate for the same statement used
$|D^\eps-D^0|\le\eps(V_{AB}+V_{BA})$, which certifies only $\eps\le1/500$ and
is exhausted near $\eps\approx.0021534$. Its exhaustion was, correctly, never
interpreted as evidence of a true sign crossing; the present theorem shows that
it was indeed an artefact of the bound, by a factor of $16.8$. Direct numerical
integration of the backward system, which is a diagnostic and not a proof,
shows no sign change anywhere in $0\le\eps\le1$ at $h=\log\frac87$.
\end{remark}

The affine form $D^0=s_0(\cp-c)$ converts Proposition~\ref{prop:reserve}
directly into a decision rule that does not depend on choosing witnesses.

\begin{corollary}[Uniform decision band]\label{cor:band}
Put $\gamma_-(x)=W_{AB}/s_0(x)$ and $\gamma_+(x)=W_{BA}/s_0(x)$. For every
pulse length, every $\eps\ge0$ and every $c\in[-\tfrac14,\tfrac14]$,
\[
 c<\cp(x)-\eps\,\gamma_-(x)\ \Longrightarrow\ AB\ \text{strictly preferred},
 \qquad
 c>\cp(x)+\eps\,\gamma_+(x)\ \Longrightarrow\ BA\ \text{strictly preferred}.
\]
At $x=7/8$, $\gamma_-\approx1.64695$ and $\gamma_+\approx1.73396$. As $h$
decreases, $\gamma_\pm/h$ decreases towards about $10.94$: the computed values
of $\gamma_+/h$ are $12.99,\ 11.21,\ 11.07,\ 10.95,\ 10.94$ at
$h=\log\frac87,\ \log\frac{50}{49},\ \log\frac{100}{99},\ \log\frac{1000}{999},\
\log\frac{10^4}{10^4-1}$.
\end{corollary}

This should be compared with the legacy cubic estimate of
\S\ref{sec:legacy}, which yields a band of half-width $16000\,h$, valid only
for $\eps\le1$ and centred at $c=0$ rather than at $\cp$. The band above is
centred correctly, shrinks linearly in $\eps$, and has half-width about
$11h\eps$ --- better by more than three orders of magnitude at $\eps=1$, and by
correspondingly more for small $\eps$.

\subsection{Short pulses: the whole range of descendant division rates}

The certified range $\eps_\star$ grows without bound as the pulse shortens,
because the contrast is of order $h^3$ while the reserve is of order $h^4$.
There is therefore a finite pulse length at which the reversal is certified for
every descendant division rate up to and including the founder's own rate.

\begin{theorem}[Short pulses]\label{thm:short}
Take $h=\log\frac{50}{49}\approx.0202027$, so that $\cp\approx.0181305$. Then
at $c=-1/4$ the schedule $AB$ has strictly smaller survival risk, and at
$c=+1/4$ the schedule $BA$ does, uniformly for \emph{every} $\eps\in[0,1]$. The
certified supremum is $\eps_\star\approx1.0238$.
\end{theorem}

\begin{proof}
Identical to Theorem~\ref{thm:pulse}, with $x=49/50$. Exact evaluation gives
$s_0\approx2.81790\times10^{-5}$, $W_{AB}\approx5.90140\times10^{-6}$ and
$W_{BA}\approx6.38199\times10^{-6}$, whence
$s_0(\cp+\tfrac14)-W_{AB}>0$ and $s_0(\tfrac14-\cp)-W_{BA}>0$; the smaller of
the two residual gaps is about $1.519\times10^{-7}$.
\end{proof}

\begin{remark}
This is the statement that the previously available small-pulse bound
(\S\ref{sec:legacy}) was reaching for, and it improves it in pulse length by
more than four orders of magnitude: the cubic remainder certifies $c=\pm1/5$ at
$\eps\le1$ only for $h\le1/80000$, whereas $\log\frac{50}{49}\approx2.02\times10^{-2}$.
Under the illustrative scale of \S\ref{sec:scale} --- a mean founder waiting
time of $24$ hours --- Theorem~\ref{thm:short} concerns pulses of about $29$
minutes each, with descendants dividing exactly as fast as the founder. That is
a qualitatively different regime from a pulse of $86$ milliseconds, and it
removes the most serious restriction of the earlier treatment.
\end{remark}

\begin{figure}[tb]
\centering\includegraphics[width=\linewidth]{figures/epsilon.pdf}
\caption{The certified descendant division rate $\eps_\star$ as a function of
pulse length, for witnesses at $c=\pm\frac14$ and $c=\pm\frac15$. Solid and
dashed blue: the certificate of Proposition~\ref{prop:reserve}. Grey: the
occupation certificate it replaces. Both families terminate on the right where
$\cp$ leaves the feasible range and the reversal genuinely ceases to exist
(Proposition~\ref{prop:boundary}); the separation between blue and grey, by
contrast, is pure slack in the older bound. The horizontal line at $\eps=1$
marks descendants dividing as fast as the founder; the blue curves cross it at
pulse lengths near $\log\frac{50}{49}$, which is Theorem~\ref{thm:short}.}
\label{fig:epsilon}
\end{figure}

\subsection{The legacy small-pulse remainder}\label{sec:legacy}

For completeness we record the uniform estimate that the above supersedes,
since it is proved by a different technique --- a semigroup remainder rather
than a differential inequality --- and its symmetric form is occasionally
convenient. For $0\le\eps\le1$ the source satisfies
\begin{equation}\label{eq:tiny}
 \bigl|F_{AB}-F_{BA}+4ch^3\bigr|\le 64000\,h^4 ,
\end{equation}
so $c>16000h$ favours $BA$ and $c<-16000h$ favours $AB$; at $c=\pm1/5$ this
requires $h\le1/80000$. Appendix~\ref{app:tiny} proves the bound. Its cubic
term comes from a founder division followed by two daughter deaths; daughter
divisions first enter at fourth order, which is why the estimate is uniform in
$\eps$ but so lossy in $h$. The shift of the true boundary away from zero is
consistent with the fourth-order term $D(0)=\tfrac72h^4+O(h^5)$.

\section{Measuring the hidden coordinate}\label{sec:measurement}

Suppose a founder division yields a paired read of both daughters, with
independent symmetric marker flips of known probability $\eta$. Coding the
states as $\pm1$ shows that the agreement indicator has success probability
\begin{equation}\label{eq:bit}
 p(c)=\tfrac12+2kc,\qquad k=(1-2\eta)^2 ,
\end{equation}
since each conditional marker mean is attenuated by $1-2\eta$ and conditional
independence multiplies the two factors. Thus $m$ independent accepted founder
divisions give $X\sim\Bin\bigl(m,p(c)\bigr)$. At $\eta=1/10$ this is
$p(c)=\tfrac12+1.28\,c$. We do not treat correlated read errors as independent
calibration noise; see Proposition~\ref{prop:corr}.

This single bit is exactly the family-level feature that
Theorem~\ref{thm:branch} removes. It does not reconstruct the hidden subtree;
it measures the one functional of the joint daughter law that the decision
depends on.

\subsection{Constant actions: a forty-division classification theorem}

At $T=\log\frac87$ and $\eps=0$, first-event integration under constant actions
gives
\begin{equation}\label{eq:constant}
 F_A=a+sc,\quad F_B=b,\qquad
 a=\tfrac{1684141}{335544320},\quad
 s=\tfrac{428807}{251658240},\quad
 b=\tfrac{6767}{1310720},
\end{equation}
obtained by integrating $G_c\bigl(1-e^{-2(T-t)},1-e^{-4(T-t)}\bigr)e^{-t}$ over
$0<t<T$, each exponential term contributing $(x-x^r)/(r-1)$ with $x=7/8$ and
constant term $1-x$. The contrast is $D_C=F_A-F_B=s\,(c-c_*)$ with
\begin{equation}\label{eq:cstar}
 c_*=\frac{b-a}{s}=\frac{144633}{1715228}\approx.0843229 ,
\end{equation}
so for constant actions \emph{high} $c$ favours $A$ --- the opposite
orientation to the pulse comparison, and a reason to keep the two thresholds
strictly separate.

\begin{theorem}[Specified separated decision task]\label{thm:forty}
Assume $\eta=1/10$, $0\le\eps\le1/4$ and $|c-c_*|\ge1/10$. From $40$
independent founder-division agreement bits, choose $A$ if $X\ge25$ and $B$
otherwise. Then the probability of choosing the better constant treatment
exceeds $.95$ throughout this class.
\end{theorem}

\begin{proof}
By Proposition~\ref{prop:reserve} applied to the two one-phase schedules
(Lemma~\ref{lem:criteria}(a) again supplies \eqref{eq:supersol}, since
$2\le4\le2\cdot2$ and $3\le3\le2\cdot3$),
$D_C^0-\eps W_A\le D_C^{\eps}\le D_C^0+\eps W_B$. On the separated class
$|D_C^0|\ge s/10=428807/2516582400\approx1.70382\times10^{-4}$, while
$\tfrac14\max(W_A,W_B)=\tfrac14W_A\approx1.50323\times10^{-4}<s/10$. Hence the
sign of $D_C^{\eps}$ agrees with the sign of $c-c_*$ throughout the class, and
classifying $c$ relative to $c_*$ determines the true recurring-source ranking.

For the classification itself, the extreme Bernoulli probabilities nearest the
decision boundary are
$p_L=\tfrac12+2k(c_*-\tfrac1{10})=\tfrac{51449691}{107201750}$ and
$p_U=\tfrac12+2k(c_*+\tfrac1{10})=\tfrac{78893339}{107201750}$. Coupling
Bernoulli trials through common uniform variables shows that binomial upper
tails increase with $p$ and lower tails decrease, so it suffices to check
\[
 \sum_{j=25}^{40}\binom{40}{j}p_L^{\,j}(1-p_L)^{40-j}<\tfrac1{20},
 \qquad
 \sum_{j=0}^{24}\binom{40}{j}p_U^{\,j}(1-p_U)^{40-j}<\tfrac1{20}.
\]
Exact rational summation proves both; the two error probabilities are
approximately $.0464141$ and $.0422845$.
\end{proof}

\begin{remark}
The admissible range $\eps\le1/4$ is $250$ times the range previously available
for this theorem, and the exact supremum is $s/(10W_A)\approx.2834$. The
decision rule and its two error probabilities are unchanged; only the
robustness statement improves, because the reserve shrank. This theorem
\emph{assumes} separation: it does not certify that assumption from the data,
does not provide a confidence interval, does not permit abstention, does not
apply near the tie, and does not use the pulse threshold.
\end{remark}

\subsection{Full-class guarantees with an unresolved output}

Without a separation assumption, the honest output has three values. An exact
equal-tailed binomial confidence interval \citep{ClopperPearson1934} for $p$
inverts through \eqref{eq:bit} to a set of admissible $c$; mapping that set
through the affine reference contrast and adding the coupling reserves to the
two ends gives an interval for the true contrast, and a declaration is made
only when that interval has one sign. Shared nuisance parameters must remain
shared when the contrast is formed.

Equivalently, and this is the implementation we certify: for a threshold $t$, a
reference slope magnitude $s_0>0$ and reserves $\rho_-,\rho_+\ge0$, define
\begin{equation}\label{eq:pends}
 p_-=\tfrac12+2k\bigl(t-\rho_-/s_0\bigr),
 \qquad
 p_+=\tfrac12+2k\bigl(t+\rho_+/s_0\bigr).
\end{equation}
When both lie in $[0,1]$, declare the low-$c$ action if
$\Prob_{p_-}\bigl(\Bin(m,p_-)\le X\bigr)\le\alpha/2$; declare the high-$c$
action if $\Prob_{p_+}\bigl(\Bin(m,p_+)\ge X\bigr)\le\alpha/2$; otherwise
return \emph{unresolved}. If an endpoint falls outside the parameter range the
corresponding decision is vacuous or forced and can be read off the whole
contrast range directly; in every example below both endpoints lie strictly
inside $[0,1]$.

\begin{proposition}\label{prop:abstain}
The above test makes an incorrect declaration with probability at most $\alpha$
throughout the admitted source class, without any separation assumption. For
known $\eta=1/10$, $\alpha=.05$ and $0\le\eps\le1/100$:
\begin{enumerate}[nosep,leftmargin=*]
\item \emph{Constant actions}, $t=c_*$, $s_0=s$, $\rho_-=\eps W_B$,
      $\rho_+=\eps W_A$, $m=200$: declare $B$ at $X\le106$ and $A$ at
      $X\ge137$. The resolution probability at $c=.2$ is $\approx.9909068$.
\item \emph{Pulses} at $h=\log\frac87$, $t=\cp$, $s_0$ as in \eqref{eq:slope},
      $\rho_-=\eps W_{AB}$, $\rho_+=\eps W_{BA}$, $m=1600$: declare $AB$ at
      $X\le1009$ and $BA$ at $X\ge1153$. The resolution probability at $c=.2$ is
      $\approx.9994762$.
\item With the same pulse rule and $m=800$: declare $AB$ at $X\le496$ and $BA$
      at $X\ge584$, with resolution $\approx.9590464$ at $c=.2$.
\end{enumerate}
All three resolution probabilities exceed $.95$ by exact rational comparison,
and at $c=-.2$ each exceeds $.999999$.
\end{proposition}

\begin{proof}
A low-$c$ declaration can be wrong only if the low-$c$ action is not in fact
better, which by \eqref{eq:contrast-band} requires the true $c$ to be at least
$t-\rho_-/s_0$ and hence $p\ge p_-$; the probability of the low-$c$ acceptance
region is then at most its value at $p_-$, which is at most $\alpha/2$ by
monotonicity of the binomial lower tail. Symmetrically for a wrong high-$c$
declaration, and the union bound gives $\alpha$. The acceptance regions are
integer tails, so their probability at any specified rational $p$ is an exact
finite sum of rationals; the cutoffs and resolution probabilities above are
computed that way, with the displayed decimals descriptive and the comparisons
against $.95$ exact integer inequalities.
\end{proof}

\begin{remark}
Three separate things improved at once here, and they should not be conflated.
The reserve shrank, which narrows the unresolved band; the admissible $\eps$
grew tenfold; and the resulting budget fell. Concretely, the previously
published rule at $\eps\le.001$ and $m=1600$ declared $AB$ at $X\le982$ and
$BA$ at $X\ge1177$, with resolution $.9721808$ at $c=.2$. The new rule at ten
times the $\eps$ range resolves more often with the same $m$, and item~3 shows
that half the sample now suffices to clear the same $.95$ resolution target.
\end{remark}

\begin{remark}
False-declaration probability is controlled \emph{unconditionally}. This is not
a claim that the conditional probability of error given a declaration is at
most $.05$. ``Resolution'' means producing either declaration; correct
resolution is a different event, although wrong declarations are bounded. Near
a robust tie an unresolved output is expected and appropriate.
\end{remark}

\subsection{Effort, and a matched direct-endpoint benchmark}

A founder divides before its exponential deadline with probability
$y=1/(1+\lambda)$, so collecting until $m$ divisions yields i.i.d.\ marks and a
negative-binomial expected root count $m/y$. The expected observed founder time
per root is $1/(1+\lambda)$, so the expected total founder time is $m$; these
are expectations, not fixed-budget guarantees. Imposing a cap $C$, with its
flag retained, changes the yield to $y_C=\bigl(1-e^{-(1+\lambda)C}\bigr)/(1+\lambda)$
while leaving the conditional bit law unchanged, because time and daughter draw
are independent. With no timeout every immortal founder eventually divides, so
exactly $m$ roots suffice and the expected total founder time is still $m$: a
timeout limits the length of an individual record but is not root-optimal in
this idealised source. Calibration reads and the paired-read itself are
additional costs, charged separately.

Direct endpoint randomisation needs no sister assay at all, and a fair
comparison must therefore fix the same task. We use the separated class of
Theorem~\ref{thm:forty}, the same $.05$ error, known nuisance bounds and no
abstention. Let $W$ be the difference of extinction indicators under $A$ and
under $B$ from two independent experimental roots. By Lemma~\ref{lem:mono},
$F_A^{\eps}\le F_A^0$ and $F_B^{\eps}\le F_B^0$, so throughout the class
$F_A,F_B\le\max(a+s/4,\,b)<3/500$ with no reserve needed at all; hence
$\Var(W)\le.012$ and $|W-\E W|\le2$. With $n$ independent pairs Bernstein's
inequality \citep[\S2.8]{BoucheronLugosiMassart2013} gives
\begin{equation}\label{eq:bernstein}
 \Prob(\text{wrong empirical sign})
 \le\exp\Bigl\{-\frac{ng^2}{2\bigl(.012+2g/3\bigr)}\Bigr\},
 \qquad g=\tfrac{s}{10}-\eps W_A .
\end{equation}
At $\eps\le1/100$ this is $g\approx1.6438\times10^{-4}$ and the right-hand side
falls below $.05$ at $n=2\,688\,960$ per arm, certified using $\log20<3$; that
is $5\,377\,920$ roots in total, against $7\,814\,272$ under the previous
reserve. It remains a sufficient, not a necessary, cost: no optimal
sample-complexity separation is proved, and Table~\ref{tab:cost} does not
assert an optimal experimental advantage. Direct outcomes retain the important
advantage of not requiring the mechanistic and marker assumptions that the
paired protocol needs.

\begin{table}[tb]
\centering\small
\caption{Different guarantees and their costs. Pair-record rows use known
$\eta=.1$; expected root counts use $\lambda=1$ without a cap. All observations
are independent roots or independently accepted founder divisions. The $\eps$
column is the certified descendant-division range for that row.}
\label{tab:cost}
\begin{tabularx}{\linewidth}{@{}lXllr@{}}
\toprule
Task & Guarantee and assumption & $\eps$ & Divisions & Roots\\\midrule
Constant classification & $.95$ correct; assumes $|c-c_*|\ge.1$ & $\le1/4$ & $40$ & $80$ exp.\\
Constant abstention & full-class error $\le.05$; resolution $.9909$ at $c=.2$ & $\le1/100$ & $200$ & $400$ exp.\\
Pulse abstention & own threshold; error $\le.05$; resolution $.9995$ at $c=.2$ & $\le1/100$ & $1600$ & $3200$ exp.\\
Pulse abstention, lean & as above; resolution $.9590$ at $c=.2$ & $\le1/100$ & $800$ & $1600$ exp.\\
Direct constant endpoint & same separated task; Bernstein sufficient bound & $\le1/100$ & --- & $5\,377\,920$\\
\bottomrule
\end{tabularx}
\end{table}

\section{Calibration, correlated errors and residual ambiguity}
\label{sec:calibration}

\subsection{What a sign needs, and what a threshold needs}

Unknown but symmetric marker error does not always prevent a decision. If
$k>0$ then \eqref{eq:bit} gives
\begin{equation}\label{eq:signread}
 \sgn\bigl(p-\tfrac12\bigr)=\sgn(c),\qquad |c|\ge\tfrac12\bigl|p-\tfrac12\bigr| ,
\end{equation}
both statements free of any knowledge of $\eta$. Combining
\eqref{eq:signread} with Corollary~\ref{cor:band} gives a genuinely
calibration-free rule at \emph{finite} pulse length, in two asymmetric halves.

\begin{proposition}[Calibration-free declarations]\label{prop:calfree}
Fix $x$ and $\eps$ with $\cp(x)-\eps\gamma_-(x)>0$, and let $I$ be a
$(1-\alpha)$ confidence interval for $p$ built without knowledge of $\eta$.
\begin{enumerate}[nosep,leftmargin=*]
\item If $I\subset[0,\tfrac12)$ then $AB$ is strictly preferred, with error
      probability at most $\alpha$.
\item If $I\subset\bigl(\tfrac12+2\bigl(\cp(x)+\eps\gamma_+(x)\bigr),1\bigr]$
      then $BA$ is strictly preferred, with error probability at most $\alpha$.
\end{enumerate}
At $x=7/8$ and $\eps\le1/100$ the hypothesis of the proposition holds, since
$\cp-\eps\gamma_-\approx.12059>0$, and the second threshold is
$\approx.80879$.
\end{proposition}

\begin{proof}
On the coverage event, in case~1 we get $p<\tfrac12$, hence $c<0$ by
\eqref{eq:signread}, hence $c<\cp-\eps\gamma_-$ by hypothesis, and
Corollary~\ref{cor:band} applies. In case~2 we get
$p-\tfrac12>2(\cp+\eps\gamma_+)$, hence
$|c|\ge\tfrac12|p-\tfrac12|>\cp+\eps\gamma_+$ with $c>0$, and
Corollary~\ref{cor:band} applies again.
\end{proof}

The asymmetry is intrinsic and worth stating plainly. Certifying $AB$ needs
only the sign of $c$, because the threshold is on the other side of zero;
certifying $BA$ needs $|c|$ to exceed a positive threshold, and a
calibration-free lower bound on $|c|$ loses a factor of two, so it demands a
strong channel. At $\eta=.1$ one has $p(1/4)=.82>.80879$, so case~2 is
feasible, but only for nearly extremal covariance. There is no uniform sample
bound as $k\to0$: at pure noise no signal remains.

The previously available calibration-free rule required
$I\subset(\tfrac12+32000h,1]$, which at $h=\log\frac87$ demands $p>4273$ and is
vacuous; it was usable only in the regime $h\le10^{-6}$. The rule above is
usable at the pulse lengths the rest of the paper works with.

Calibration is genuinely needed to locate $c$ against a \emph{shifted}
threshold. For example $(c,\eta)=(.05,0)$ and $(.2,.25)$ both give $p=.6$, yet
lie on opposite sides of $c_*$; no amount of agreement data resolves that pair.
A calibration interval must then be projected jointly with the bit confidence
set, and the union-bound coverage event for the two measurements is sufficient.
Estimating the channel does not turn it into a known constant.

\subsection{Asymmetric and correlated read channels}

For a common asymmetric binary channel with $\E[Y_i\mid Z_i]=a+bZ_i$,
$\Prob(Z_i=1)=\tfrac12$, $\Cov(Z_1,Z_2)=c$ and conditionally independent read
errors,
\begin{equation}\label{eq:asym}
 \Cov(Y_1,Y_2)=a^2+ab+b^2(\tfrac14+c)-(a+\tfrac b2)^2=b^2c .
\end{equation}
The centred observed covariance therefore preserves the sign of $c$ whenever
$b\ne0$. This requires the marginal read frequencies as well as the paired
reads: an agreement bit alone does not identify the unknown centre.

Correlated read errors are a more serious threat, and the right response is to
price them rather than to assume them away.

\begin{proposition}[Additive reserve for correlated read errors]\label{prop:corr}
Suppose $\E[Y_i\mid Z_1,Z_2]=a+bZ_i$, so that the conditional mean of one read
does not itself acquire a dependence on the sister state, and suppose the two
read errors have conditional covariance $q(Z_1,Z_2)$. Then
\begin{equation}\label{eq:corr}
 \Cov(Y_1,Y_2)=b^2c+\E\bigl[q(Z_1,Z_2)\bigr].
\end{equation}
Consequently, if $\bigl|\E q\bigr|\le\delta_{\mathrm{err}}$ is known, then a
confidence interval for the observed covariance lying entirely above
$\delta_{\mathrm{err}}$ certifies $c>0$, and one lying entirely below
$-\delta_{\mathrm{err}}$ certifies $c<0$.
\end{proposition}

\begin{proof}
By the law of total covariance,
$\Cov(Y_1,Y_2)=\Cov\bigl(\E[Y_1\mid Z],\E[Y_2\mid Z]\bigr)
+\E\bigl[\Cov(Y_1,Y_2\mid Z)\bigr]$. The first term is
$\Cov(a+bZ_1,a+bZ_2)=b^2c$ and the second is $\E[q]$.
\end{proof}

The value of \eqref{eq:corr} is not that it rescues the assay --- it does not
--- but that it shows precisely what must be measured: a bound on the
\emph{average} error covariance, not a claim that errors are independent. Note
that a certified sign of $c$ combines with Proposition~\ref{prop:calfree}
case~1 to certify $AB$; certifying $BA$ still needs a magnitude.

\subsection{Irreducible ambiguity, and decisions under it}

On the two sources of Theorem~\ref{thm:pulse} the retained-record laws are
identical, so any possibly randomised binary selector has worst-case error at
least $1/2$, for any number of independent records: if $E$ denotes the event of
choosing one of the actions, the two error probabilities sum to
$1-P_1(E)+P_2(E)=1$. With unequal data laws the same two-point argument gives
$\max\Prob(\text{error})\ge\bigl(1-\TV(P_1,P_2)\bigr)/2$, and multiplying by a
common gap lower bound converts this into an excess-risk lower bound. For the
paired assay the one-root information is $y\,\kl(p_1\Vert p_2)$ and
independent-root Kullback--Leibler divergences add. These are standard testing
tools \citep{Tsybakov2009,KaufmannCappeGarivier2016} specialised to this record
law, not new information inequalities.

When the remaining gap makes a ranking costly, report the uncertainty in risk
units. If an estimated contrast has absolute error at most $e$, then a wrong
sign costs at most $e$ in excess risk; $2e$ is the correct bound when the two
risks were estimated separately each with error $e$, and the two conventions
must be kept distinct. On a contrast interval $[L,U]$, choosing the first
action has worst regret $\max(U,0)$ and the second $\max(-L,0)$; for $L<0<U$,
randomising with probability $-L/(U-L)$ on the first action balances the two
endpoint losses and minimises the worst regret at
\begin{equation}\label{eq:regret}
 \Gamma=\frac{(-L)\,U}{U-L}
\end{equation}
\citep{Stoye2012,Manski2004}. A declared tolerance $\delta$ can therefore
justify an action even when the ranking is unresolved, and coverage failure
adds at most $\alpha$ to the unconditional expected regret.

\section{A controlled refinement, and what the numbers mean}
\label{sec:refinement}

\subsection{Mother-dependent phenotype marginals}

The source \eqref{eq:kernel} deliberately separates sister coupling from
marginal phenotype persistence, so that the theorem cannot be explained away as
an artefact of inheritance. To include genuine persistence, keep the founder
law \eqref{eq:kernel} but give a type-$i$ descendant the pair law
\begin{equation}\label{eq:persistent}
 K_i=\begin{pmatrix}
 p_i^2+\kappa_i & p_i(1-p_i)-\kappa_i\\
 p_i(1-p_i)-\kappa_i & (1-p_i)^2+\kappa_i\end{pmatrix},
 \qquad p_S=.6,\ p_T=.4,\ \kappa_i\in\{-.1,.1\},
\end{equation}
whose entries are all nonnegative (they are, in the four cases,
$\tfrac{13}{50},\tfrac{17}{50},\tfrac{17}{50},\tfrac3{50}$ and permutations).
Now a daughter matches her mother's type with probability $.6$. The founder
covariance $c$ and the descendant covariances $\kappa_i$ are separate
parameters and must be named separately.

\begin{corollary}\label{cor:memory}
Theorems~\ref{thm:pulse} and \ref{thm:short}, with the same $\eps$ ranges and
the same gaps, remain valid when descendants use \eqref{eq:persistent}. More
generally they remain valid when descendants use \emph{any} type-dependent
offspring law on $\{S,T\}$ whatsoever. Sources sharing the founder's daughter
marginals and the event rates, but differing in their founder or descendant
sister covariances, remain observationally equivalent under the retained-branch
protocol.
\end{corollary}

\begin{proof}
The descendant offspring law enters the backward system \eqref{eq:field} only
through the terms $\eps\,(G-\cdot)$ in the two daughter coordinates, and
Lemma~\ref{lem:mono} and Proposition~\ref{prop:reserve} were stated for an
arbitrary such $G$. The only hypothesis that refers to $G$ is
\eqref{eq:supersol}, and Lemma~\ref{lem:criteria}(b) supplies it for every
kernel once $x\ge x_\flat$. Both theorems use $x=7/8$ and $x=49/50$, which
exceed $x_\flat$, and indeed every pulse length at which the reversal exists
does. The reference contrast at $\eps=0$ and the reserves $W_a$ are computed
from the $\eps=0$ trajectory alone and are therefore untouched, so the
certificates transfer verbatim. The source remains nonexplosive. Finally
Theorem~\ref{thm:branch} uses only the common daughter marginals, which
\eqref{eq:persistent} shares across $\kappa_i$, giving the observation claim.
\end{proof}

This is a robustness theorem, not a calibration of real inherited phenotypes.
Its content is that sister dependence and mother-to-daughter memory are
distinct model features and that only the former is at stake here: the
descendant mechanism may be arbitrary, and the reversal is unaffected.

\subsection{Physical scale}\label{sec:scale}

Let the physical founder division hazard be $\beta$. Model time is
$\beta t_{\mathrm{phys}}$; the descendant division hazard is $\eps\beta$, the
death hazards are the model values times $\beta$, and a pulse lasts $h/\beta$.
Rescaling time changes all dimensional rates together and therefore cannot
repair their ratios. Note also that an exponential mean division waiting time
is not a population doubling time, and that the reciprocal descendant hazard is
a latent division-clock scale, not the realised waiting time of a cell that can
die first.

\begin{table}[tb]
\centering\small
\caption{Model quantities under the illustrative mean founder waiting time
$1/\beta=24$ hours. The last four rows are the point of \S\ref{sec:recurring}:
the certified regime has moved from a latent descendant division clock of a
thousand days to one of a few weeks, and --- at a half-hour pulse --- to no
restriction at all.}
\label{tab:scale}
\begin{tabularx}{\linewidth}{@{}llX@{}}
\toprule
Quantity & Physical reading & Status\\\midrule
$h=\log\frac87$ & $3.20$ h per pulse, $6.41$ h in total & worked example\\
$h=\log\frac{50}{49}$ & $29.1$ min per pulse, $58.2$ min in total & Theorem~\ref{thm:short}\\
$h=10^{-6}$ & $86.4$ ms per pulse & legacy bound \eqref{eq:tiny}\\
$\eps=.001$ & descendant division clock $24\,000$ h $=1000$ d & previously certified\\
$\eps=1/30$ & descendant division clock $720$ h $=30$ d & Theorem~\ref{thm:pulse}\\
$\eps=1/16$ & descendant division clock $384$ h $=16$ d & Theorem~\ref{thm:pulse}\\
$\eps=1$ & descendants divide as fast as the founder & Theorem~\ref{thm:short}\\
\bottomrule
\end{tabularx}
\end{table}

The earlier restriction to $\eps\le.001$ was a genuine weakness: a population
whose founders and daughters have comparable division rates was simply outside
the certified regime. Theorem~\ref{thm:short} removes that restriction
entirely, at the cost of a shorter pulse, and Theorem~\ref{thm:pulse} relaxes
it by more than an order of magnitude at the original pulse length. What
remains restrictive is the specification of the response law \eqref{eq:rates}
and the instantaneous nonperturbing readout, not the demographic ratio.

\subsection{Effect sizes and multiple founders}

At $c=.2$ and $\eps=0$,
\[
 F_{AB}=\tfrac{40043373837487}{1319413953331200}\approx.030349364,
 \qquad
 F_{BA}\approx.030567813 .
\]
The absolute survival-risk improvement is $\approx2.18449\times10^{-4}$, that
is $.0218$ percentage points, while the relative increase in one-founder
extinction is about $.72\%$. These are different effect-size descriptions and
should not be interchanged. Both schedules leave survival near $97\%$, so the
example establishes a ranking mechanism, not effective eradication.

For $n$ independent $P$ founders, extinction is $F_a^n$ and
\begin{equation}\label{eq:founders}
 \bigl|F_{AB}^n-F_{BA}^n\bigr|
 \le n\max(F_{AB},F_{BA})^{n-1}\bigl|F_{AB}-F_{BA}\bigr| .
\end{equation}
The ranking is preserved for every $n$, but the magnitude collapses: at $n=2$
the extinction difference is $\approx1.331\times10^{-5}$ and at $n=10$ it is
$\approx4.930\times10^{-17}$, the individual extinction probabilities there
being $\approx6.630\times10^{-16}$ and $\approx7.123\times10^{-16}$. Random
preparations require averaging the full extinction product rather than
substituting a mean count, and a shared environment requires conditional
analysis before any family factorisation.

\section{Scope and interpretation}\label{sec:scope}

\subsection{What the obstruction is specific to}

The obstruction is specific to a pruned observation and a family-level
endpoint. Exact full-state molecular conservation removes it: if a mother
contains $(a,r)$ molecules and the observed daughter contains $(i,j)$, the
sister must contain $(a-i,r-j)$ under pure partitioning, so the joint law is
determined by the one-daughter law. Our phenotype covariance family must not be
treated as a free molecular covariance under those stronger constraints.

Founder immortality before division, exponential clocks, known intervention
responses, instantaneous nonperturbing reads, conditionally independent marker
errors and independent families are all substantive assumptions. Persistent
marginals are covered by Corollary~\ref{cor:memory}; arbitrary reporter delay,
unobserved age structure, unknown responses and shared environments are not.
Correlated reads are covered only to the extent that
Proposition~\ref{prop:corr} is supplied with a defensible $\delta_{\mathrm{err}}$.

The endpoint \eqref{eq:risk} uses zero terminal extinction data. Eventual
extinction after withdrawal is governed by the extinction vector of the
withdrawal source and may erase these finite-time differences entirely.
Clinical cure, relapse, detection and resistance establishment are different
events. Nor does comparing this finite menu identify a globally optimal feedback
policy.

\subsection{Before invoking Theorem~\ref{thm:branch} for a real protocol}

Theorem~\ref{thm:branch} is a statement about one precisely specified
observation rule. A platform may resemble it superficially and fail its
hypotheses. Table~\ref{tab:platform} lists the checks we consider mandatory.
None of the entries is an assertion that a named platform does or does not
satisfy them; establishing that is a separate, empirical task that this paper
does not perform.

\begin{table}[tb]
\centering\small
\caption{Protocol correspondence checks. A nonuniform but state-independent
retention rule would admit an analogous theorem provided the induced daughter
marginals still agree; a physically selected branch requires its own proof.}
\label{tab:platform}
\begin{tabularx}{\linewidth}{@{}lX@{}}
\toprule
Candidate data situation & Required check\\\midrule
Microfluidic devices that follow one branch \citep{Wang2010}
 & Is retention independent of state, or do pole age, position or survival
   select the branch? Include the relevant state variables and derive the
   actual kernel.\\
Destructive single-cell sampling
 & This is generally a snapshot protocol, not a complete retained history.
   Establish its own observation law; do not apply Theorem~\ref{thm:branch}
   verbatim.\\
Subsampled retrospective lineage reconstruction
 & Does the record retain shared ancestry, sister pairs or family sizes that
   reveal dependence? Derive the sampling law.\\
Full time-lapse lineage trees \citep{Sandler2015,Spencer2009}
 & Both sisters may be observed, so the exact pruning obstruction can vanish.
   Investigate the actual resolution and error model --- this is the regime in
   which the repair of \S\ref{sec:measurement} is natural.\\
\bottomrule
\end{tabularx}
\end{table}

\subsection{The narrow conclusion}

The mathematical conclusion is narrow and we state it without ornament: a
specified missing family-level measurement can reverse a specified schedule
comparison, this remains true when descendants divide at any rate up to the
founder's own, and the uncertainty relevant to that comparison can be bounded
without recovering every parameter. Observing how an individual cell behaves and
choosing an intervention for its family are different inference problems, and
the right measurement is dictated by the continuation products that appear in
the decision. A small paired observation can resolve a comparison that complete
retained-branch histories cannot resolve at all.

\appendix

\section{Exact certificates}\label{app:cert}

\paragraph{The boundary.} The polynomial appearing in the derivative proof of
Proposition~\ref{prop:boundary} is \eqref{eq:P}; all eleven coefficients are
positive. Expanding at $x=1$ gives $\mathcal A(1)=210$ and $\mathcal Q(1)=735$.
The equation $\cp=1/4$ is, after clearing denominators,
\[
 -6x^9-18x^8-36x^7-60x^6-76x^5-78x^4-66x^3-13x^2+81x+62=0 ,
\]
and rational substitution at $x=.80120$ and $x=.80121$ brackets its unique root
in $(0,1)$; monotonicity of $\cp$, not a floating-point root count, proves the
uniqueness.

\paragraph{Reserves and margins.} At $h=\log\frac87$ the pulse threshold is
$8967219299/65423564404$ and the affine slope magnitude is
$114491237707/32985348833280$. The reserve polynomials of
Table~\ref{tab:reserves} are evaluated using a rational bracket for
$\log\frac87$ of width below $10^{-30}$, obtained from the alternating series
for $\log(1+y)$ at $y=1/7$; for the coarser statements the bound
$\log\frac87<134657/1008420$ suffices. The residual gaps quoted in
Theorem~\ref{thm:pulse} are
\[
 \min\Bigl(D(-\tfrac15)-\tfrac1{30}W_{AB},\ -D(\tfrac15)-\tfrac1{30}W_{BA}\Bigr)
 \approx1.78315\times10^{-5},
\]
and the corresponding quantity at $\eps=1/16$, $c=\mp1/4$ is
$\approx1.58398\times10^{-5}$. For the record, the previously published
occupation certificate at $\eps=1/500$ left the gap
$154033465997/9895604649984000>1.55\times10^{-5}$, and at $\eps=.001$ the gap
$1157858798221/9895604649984000\approx1.17007\times10^{-4}$.

\paragraph{Binomial arithmetic.} The exact binomial implementation represents
$p=a/d$ and the $j$th term by the integer $\binom mj a^j(d-a)^{m-j}$ over the
common denominator $d^m$; prefix and suffix sums then give every tail
comparison without floating-point rounding, and the same routine evaluates the
resolution probabilities of Proposition~\ref{prop:abstain}. The displayed
decimals are descriptive; the comparisons against $.95$ and $\alpha/2$ are
exact integer inequalities.

\section{The sensitivity reserve in closed form}\label{app:reserve}

The quantities of Definition~\ref{def:reserves} are elementary. On a phase with
constant death rate $r$ over $[0,u]$ measured from the deadline,
\[
 L(u)=\frac{1-e^{-ru}}{r},\qquad
 \sigma(u)=\frac{e^{ru}-1-ru}{r}\,e^{-ru}=\frac{1-e^{-ru}-rue^{-ru}}{r} .
\]
For a two-phase schedule one splits the inner integral in \eqref{eq:sigmaL} at
the phase boundary, so that $\sigma_i$ is again a combination of $1$, $u$ and
exponentials $e^{-ku}$ with integer $k$; the outer integral in \eqref{eq:VW}
then has an integrand of the form $\sum_k(\alpha_k+\beta_ku)e^{-ku}$, and each
term integrates in closed form. With $x=e^{-h}$ rational the result lies in
$\mathbb Q+\mathbb Q\,h$, the linear term in $h$ arising exactly from the
constant ($k=0$) contributions. The occupation reserve $V_a$ has no $h$ term,
since its integrand contains no polynomial factor; this is why $V_a$ is a plain
rational in Table~\ref{tab:reserves}.

Writing out the two-phase occupation reserve for one daughter type with phase
hazards $(r,s)$ gives the compact expression
\[
 V(r,s)=\frac{(1-x)-(1-x^r)/r}{r-1}
 +\frac{x-x^r}{r-1}\cdot\frac{1-x^s}{s}
 +x\,\frac{(1-x)-(1-x^s)/s}{s-1},
\]
whose three terms integrate the occupation before the switch, the survivors
present at the switch, and the daughters born after it; summing over the two
daughter types reproduces the $V_a$ column of Table~\ref{tab:reserves}. The
corresponding expression for $W_a$ is longer but is produced by the same
mechanical procedure, and the accompanying script emits it as an exact element
of $\mathbb Q+\mathbb Q\,h$.

Two sanity checks are worth recording. First, the $V_a$ column reproduces
exactly the constants obtained previously by a completely different route --- a
first-added-event coupling with the bound ``probability that an extra division
occurs'' --- which confirms that $V_a$ is the reserve that argument produces
and that $W_a$ is a strict refinement of it rather than a different quantity.
Second, direct high-accuracy numerical integration of the backward system at
$x=7/8$ gives $\lim_{\eps\downarrow0}(F^0_a-F^{\eps}_a)/\eps\approx.0035$ for
both schedules, against $W_a\approx.0057$ and $V_a\approx.0507$: the sensitivity
reserve is within a factor of $1.7$ of the true linear response, while the
occupation reserve is off by a factor of $14$.

\section{The legacy small-pulse remainder}\label{app:tiny}

We prove \eqref{eq:tiny}. Write $K(u,v)=\sum_{jk}K_c(j,k)u_jv_k$. Daughter
extinction starts as $d_at+O(t^2)$. Reverse propagation for chronological $AB$
begins with $B$; integrating the founder's offspring term over the next
backward phase gives the cubic coefficient
$K(d_A,d_A)/3+4K(d_B,d_B)/3+K(d_A,d_B)$. Interchanging $A$ and $B$ leaves the
difference $K(d_B,d_B)-K(d_A,d_A)=-4c$, and the remaining terms first contribute
at order four.

To bound the remainder, use the coefficient $\ell^1$ norm on polynomials in
$(p,s,t)$. For $\eps\le1$ every component of \eqref{eq:field} has norm at most
$10$, so if $L_af=\nabla f\cdot\Phi_a^\eps$ and $\deg f\le d$ then
$\|L_af\|_1\le10d\|f\|_1$. The first word $L_ap=G_c-p$ has norm $2$, so every
four-letter word has norm at most $2\cdot20\cdot30\cdot40=48000$. Let
$T_a(t)f=f\circ F_a^t$; invariance of the cube makes this a sup-norm
contraction, and $T_a(t)f=f+\int_0^tT_a(u)L_af\,du$. Expanding
$T_B(h)T_A(h)p$ to total degree three leaves five remainder terms with
coefficients $h^4/\bigl(i!\,(4-i)!\bigr)$, $i=0,\dots,4$, multiplying words of
length four evaluated on the cube; their coefficient sum is $2^4h^4/4!$. Each
schedule remainder is therefore at most $32000h^4$, and subtracting the two
gives \eqref{eq:tiny}.

This bound proves a genuinely uniform statement in $\eps$, which is why it was
worth having, but it is extremely lossy in $h$ and it is centred at $c=0$
rather than at the true boundary $\cp$. Both defects are removed by
Corollary~\ref{cor:band} and Theorem~\ref{thm:short}, which are uniform in
$\eps$ as well.

\section{Reproducibility and formal scope}\label{app:repro}

\paragraph{What is proved how.} The source construction, the identification of
the backward system with the extinction probabilities, the retained-history
equality of Theorem~\ref{thm:branch}, the analytic integration of
\S\ref{sec:boundary}, the comparison and differential-inequality arguments of
\S\ref{sec:recurring}, the semigroup remainder of Appendix~\ref{app:tiny} and
the concentration and testing arguments of \S\ref{sec:measurement} and
\S\ref{sec:calibration} are conventional mathematical proofs, given in full in
this paper. Every numerical constant that enters a proof step --- the factored
contrast \eqref{eq:D}, the positivity of the coefficients of $P$, the bracket
for $x_0$, the four reserve polynomials, the residual gaps, the binomial tails
and resolution probabilities, the Bernstein budget, and the feasibility of
\eqref{eq:persistent} --- is recomputed in exact rational or symbolic
arithmetic by the accompanying script, which performs $82$ assertions and
reports them individually. Where an irrational quantity is unavoidable, namely
$h=\log\frac87$ and $\log\frac{50}{49}$, it is enclosed in a rational interval
of certified width and the enclosure is propagated. No floating-point value
supports any claim in this paper; floating-point numbers appear only as
descriptive decimals.

\paragraph{What is not claimed.} We make no claim of machine-checked
verification of the stochastic content. Any formalisation accompanying earlier
versions of this material covered finite algebraic identities and rational
margins only; a description such as ``formally verified stochastic treatment
theorem'' would overstate it, and we do not use one. The numerical ODE
integrations reported in Appendix~\ref{app:reserve} are diagnostics that
motivated and cross-check the analytic bounds; no proof depends on them.

\paragraph{Reproduction.} The submission contains \texttt{main.tex},
\texttt{refs.bib}, the three vector figures, and \texttt{scripts/} holding the
exact-arithmetic engine, the certificate script and the figure generator. The
certificate script prints a labelled dump of every constant quoted above,
followed by its assertion log. It has no dependency beyond a standard Python
distribution with \texttt{sympy}; the figure script additionally uses
\texttt{numpy} and \texttt{matplotlib}, and the numerical cross-check in
Appendix~\ref{app:reserve} uses \texttt{scipy}.

\paragraph{Provenance.} A shorter and unpublished working note covering part of
this material, together with an accompanying supplement, circulated privately in
September 2026; the present paper supersedes both, and no argument here depends
on either. The exact witnesses at $c=\pm1/5$ and $h=\log\frac87$, the factored
contrast \eqref{eq:D}, the retained-branch equivalence and the forty-division
rule originate in that note; the monotonicity lemma, the sensitivity reserve,
the enlarged $\eps$ ranges, the uniform decision band, the short-pulse theorem,
the calibration-free rule at finite pulse length and the correlated-error
reserve are new here. No clinical parameters were estimated and no new data were
collected.

\bibliographystyle{plainnat}
\bibliography{refs}

\end{document}
